\documentclass[11pt,a4paper]{article}

\usepackage{amsmath}
\usepackage{mathtools}
\usepackage{fontspec}
\usepackage{unicode-math}
\setmathfont{KpMath-Regular.otf}
\setmathfont{KpMath-Bold.otf}[version=bold]
\usepackage{microtype}
\usepackage[dvipsnames]{xcolor}
\usepackage[margin=24mm, top=28mm, bottom=28mm]{geometry}

\usepackage{fancyhdr}
\pagestyle{fancy}
\fancyhf{}
\fancyhead[L]{\small\itshape Where the FFT Sphere Solver Runs}
\fancyhead[R]{\small\itshape 31 October 2026}
\fancyfoot[C]{\thepage}
\renewcommand{\headrulewidth}{0.4pt}

\setcounter{secnumdepth}{3}
\usepackage{booktabs}
\usepackage{array}
\usepackage[font=small, labelfont=bf, labelsep=period]{caption}
\usepackage{setspace}
\setstretch{1.12}
\setlength{\parskip}{0.4em}
\usepackage[colorlinks=true, linkcolor=blue!60!black, urlcolor=blue!60!black]{hyperref}
\usepackage{xurl}

\newcommand{\nsub}{n_{\mathrm{sub}}}
\newcommand{\nP}{n_{P}}
\newcommand{\code}[1]{\nolinkurl{#1}}

\title{\bfseries Where the FFT Sphere Solver Runs\\[4pt]
\large Memory, precision and hardware for the Foldy--Lax sphere,\\
measured on Graviton CPUs and an NVIDIA A10G}
\author{}
\date{31 October 2026}

\begin{document}
\maketitle

\begin{abstract}
\noindent
The FFT-accelerated Foldy--Lax solver for a voxelised sphere
(\code{compute_sphere_foldy_lax_fft_gpu}) was measured from $\nsub = 12$ to $96$ cells across the diameter,
up to 463\,400 cells and 4.2 million unknowns, on eight Graviton3 cores and on an NVIDIA A10G, at $k_S a = 1$. Four
results follow.

\emph{Memory.} The GMRES solver allocated its whole Krylov basis, $(\mathrm{maxiter}+1)$ vectors, before the first
step: 13.4\,GB of a 27.7\,GB peak at $\nsub = 96$, for solves that took four steps. It now grows the basis a vector
a step, with results bit-identical to before. What remains is dominated by the padded grids of the FFT product.

\emph{Speed.} At $\nsub = 96$ the A10G forms the FFT product 15 times faster than the eight cores in complex128, and
about 50 times faster in complex64. The solve then stops being the cost: the kernel and the right-hand sides, built on
the CPU, take most of the time.

\emph{Precision.} With the kernel built in complex128 and only the solve in complex64, the far field agrees with the
complex128 solution to $3.9\times10^{-9}$ at a GMRES tolerance of $10^{-6}$, against a discretisation error of
$6.7\times10^{-5}$. Single precision costs nothing measurable.

\emph{Design.} The two halves want different machines: the kernel build, in complex128, parallel and mostly
independent of frequency, on many CPU cores; the solve, in complex64 and limited by memory bandwidth, on a GPU. The
canonical kernel that passes between them is 160\,MB at $\nsub = 96$ in complex64.
\end{abstract}

\tableofcontents

%=====================================================================
\section{The solver and the measurement}\label{sec:setup}

\subsection{What is solved}

A sphere of radius $a = 0.5$\,m is voxelised on a cubic grid of $\nsub$ cells across its diameter; the cells whose
centres lie inside it, $N$ of them, each carry nine unknowns (displacement and Voigt strain). The Foldy--Lax system
couples every pair of cells through the cell-averaged propagator. Its product with a vector is a convolution, formed
by FFT on a zero-padded grid of $\nP^3$ points per component, with $\nP$ the smallest odd length at least
$2\nsub - 1$ whose only prime factors are 3, 5 and 7 (\code{fast_odd_size}): $\nP = 225$ at $\nsub = 96$. The
kernel is stored only at the canonical wavenumbers of the cube group, about $1/46$ of the grid
(\S\ref{sec:design}). GMRES (\code{torch_gmres}) solves one column for each orbit of the right-hand sides under the
symmetries: three Taylor patterns, which give the composite $9\times9$ T-matrix, and six plane-wave columns, which give
the far field. The other columns follow by symmetry.

The kernel is always built on the CPU in complex128: the cell-averaged pair blocks come from a wavenumber series
whose recursions need the precision (\code{cell_average_kseries.py}). Only the solve runs in the device's dtype,
on the kernel rounded to it.

\subsection{The test case}

The background is the project's test medium, $\alpha = 5000$\,m/s, $\beta = 3000$\,m/s, $\rho = 2500$\,kg/m$^3$; the
sphere's contrast is $\Delta\lambda = 2$\,GPa, $\Delta\mu = 1$\,GPa, $\Delta\rho = 100$\,kg/m$^3$; the frequency is
fixed at $k_S a = 1$, so that refining $\nsub$ refines the discretisation of the same physical problem. GMRES
converged in three or four steps at every size.

\subsection{The probe}

\code{scripts/probe_fft_sphere_scaling.py} runs each size in a fresh process, so that its peak memory is its own,
and records the time to build the kernel, each GMRES solve (iterations, products, time), the total, the peak resident
memory of the process, and on CUDA the allocator's peak. The k-series store, the universal coefficients of the pair
blocks (\code{numerics.yml}: \code{kseries_store}), grows on first use by the offsets a size needs; every run reported
here found it already grown. \code{scripts/measure_fft_sphere_precision.py} compares precisions and the exact Mie
solution (\S\ref{sec:precision}).

\subsection{The machines}

\begin{table}[h]
\centering
\caption{The two AWS instances (Sydney, \code{ap-southeast-2}), each used through the project's image
(\code{aws/README.md}): Python 3.12, NumPy 2.4.2, torch 2.10.0 (the CUDA 12.8 build on the GPU).}
\label{tab:machines}
\small
\begin{tabular}{llll}
\toprule
Instance & CPU & Memory & GPU \\
\midrule
r7g.2xlarge & 8 Graviton3 cores (arm64) & 64\,GiB & --- \\
g5.xlarge & 4 vCPUs: 2 AMD EPYC cores, hyperthreaded & 16\,GiB & NVIDIA A10G, 24\,GB \\
\bottomrule
\end{tabular}
\end{table}

The g5.xlarge was chosen for its GPU; its two physical cores matter for the stages that stay on the CPU
(\S\ref{sec:profile}). Two GPU runs that shared the instance with another job are excluded from the data; the GPU
rows here each had the instance to themselves.

%=====================================================================
\section{Memory}\label{sec:memory}

\subsection{The Krylov basis was allocated in full}

\code{torch_gmres} has no restart, and it allocated the basis for every iteration it might take,
$(\mathrm{maxiter}+1) \times 9N$ complex numbers, before the first. At the default $\mathrm{maxiter} = 200$ that is
13.4\,GB at $\nsub = 96$ in complex128, for solves that converged in four steps. Lowering $\mathrm{maxiter}$ to 20
removed almost exactly that much (table~\ref{tab:krylov}).

\begin{table}[h]
\centering
\caption{Peak resident memory on r7g.2xlarge, complex128, with the basis allocated in full. The drop is the
basis's: 3.44 against 3.56\,GB at $\nsub = 64$, 11.16 against 12.01\,GB at $\nsub = 96$.}
\label{tab:krylov}
\begin{tabular}{rrrrrr}
\toprule
& & \multicolumn{2}{c}{$\mathrm{maxiter} = 200$} & \multicolumn{2}{c}{$\mathrm{maxiter} = 20$} \\
\cmidrule(lr){3-4}\cmidrule(lr){5-6}
$\nsub$ & cells & basis (GB) & peak (GB) & basis (GB) & peak (GB) \\
\midrule
64 & 137\,376 & 3.98 & 8.88 & 0.42 & 5.44 \\
96 & 463\,400 & 13.41 & 27.71 & 1.40 & 16.55 \\
\bottomrule
\end{tabular}
\end{table}

The basis is now a list to which Arnoldi appends each vector it makes, so its memory follows the iterations taken,
$(n_{\mathrm{iter}}+1)\times 9N$, and $\mathrm{maxiter}$ can be set generously against hard cases at no cost; only the
$(\mathrm{maxiter}+1)\times\mathrm{maxiter}$ Hessenberg matrix is sized by it. The arithmetic is unchanged: against the
previous version, solutions, iteration counts and residuals are bit-identical in both precisions, for converging
solves, solves that reach $\mathrm{maxiter}$, and the breakdown branch. With $10^6$ unknowns, $\mathrm{maxiter} = 200$
and four steps, the peak fell from 3.29 to 0.24\,GB. Every GPU measurement below uses the new version.

\subsection{What remains}

Without the full basis, the A10G's peak at $\nsub = 96$ in complex128 is 14.6\,GB, of which the basis is 0.33\,GB and
the canonical kernel 0.32\,GB. One padded grid stack, nine components of $\nP^3$ complex numbers, is 1.64\,GB, so the
remainder is about eight to nine stacks: the temporaries of the FFT product (the packed grid, its transform, the
product with the kernel and its inverse transform), as inferred from their size; it has not been profiled. Two
consequences follow. Device memory scales as $\nP^3$, not as $N$, and the padding to a 3-5-7-smooth length makes
$\nP^3$ at $\nsub = 96$ $(225/191)^3 = 1.63$ times the minimum. Reducing the temporaries, by transforming in place or
component by component, is the next memory saving.

The host's peak also includes about 2.5 to 3\,GB that does not grow with $N$: 2.7 to 3.1\,GB at $\nsub = 24$, where
the solver's own arrays are well under 1\,GB. It is absent from the runs made while the k-series store held only the
offsets of small sizes (0.36\,GB at $\nsub = 12$), so the memory-mapped store is the likely part of it.

%=====================================================================
\section{Time on the CPU and on the GPU}\label{sec:time}

\begin{table}[h]
\centering
\caption{Time per FFT product, and the total for one sphere (both the composite T-matrix and the far-field columns),
in seconds. CPU: r7g.2xlarge, complex128, gmres\_tol $10^{-8}$. GPU: g5.xlarge, complex128 at $10^{-8}$ (four
steps a solve) and complex64 at $10^{-6}$ (three).}
\label{tab:time}
\begin{tabular}{rr rrr rrr}
\toprule
& & \multicolumn{3}{c}{seconds per product} & \multicolumn{3}{c}{total seconds} \\
\cmidrule(lr){3-5}\cmidrule(lr){6-8}
$\nsub$ & cells & CPU & GPU c128 & GPU c64 & CPU & GPU c128 & GPU c64 \\
\midrule
24 & 7\,208 & 0.0344 & 0.0184 & 0.0118 & 3.0 & 2.6 & 2.0 \\
32 & 17\,256 & 0.0819 & 0.0119 & 0.0074 & 6.0 & 2.7 & 2.4 \\
48 & 57\,856 & 0.437 & 0.0369 & 0.0138 & 26.9 & 6.6 & 5.2 \\
64 & 137\,376 & 1.021 & 0.0694 & 0.0237 & 60.8 & 13.2 & 10.5 \\
96 & 463\,400 & 4.911 & 0.318 & 0.096 & 283.1 & 50.2 & 37.3 \\
\bottomrule
\end{tabular}
\end{table}

At $\nsub = 96$ the A10G forms a product 15 times faster than the eight Graviton cores in complex128 (0.318 against
4.6 to 4.9\,s), and 48 to 51 times faster in complex64 (0.096\,s). Complex128 is the A10G's weak mode for arithmetic,
yet it still wins by an order of magnitude: the product is a sequence of FFTs and pointwise multiplications, limited
by memory bandwidth more than by floating-point rate. Halving the bytes, complex64 makes the product 3.3 times faster
still, and the tolerance it can reach, $10^{-6}$, takes three steps instead of four (42 products a sphere against 53
or 54).

The totals tell a different story from the products. In complex64 at $\nsub = 96$ the solves take 4.0\,s of 37.3: the
rest is on the CPU. On the r7g.2xlarge, whose eight cores ran close to 800\% busy, the kernel took 11.7\,s; on the
g5.xlarge's two cores it takes 18.5\,s; and about 15\,s more goes to stages that the probe did not separate. Section~\ref{sec:profile} finds
them.

%=====================================================================
\section{Precision}\label{sec:precision}

Each sphere was solved three times on the A10G with the same complex128 kernel: in complex128 at gmres\_tol
$10^{-10}$, the reference, and in complex64 at $10^{-5}$ and at $10^{-6}$ (single precision cannot reach the
reference's tolerance). Two errors are compared, both relative, in the maximum over the components and angles: the
\emph{precision} error, complex64 against complex128, on the composite T-matrix and on the far field; and the
\emph{discretisation} error, complex128 against the exact elastic Mie solution, on the far field. The far field is
for a $P$ wave incident along $z$, observed at nine angles in the $z$--$x$ plane at $5\times10^4$ radii, the distance at
which \code{gate_sphere_cell_average_vs_mie.py} established that the asymptotic and exact fields agree. The
discretisation error is given raw and with the voxel volume scaled to the sphere's, which removes the staircase's
dilution but not its shape.

\begin{table}[h]
\centering
\caption{The error of solving in complex64, beside the error of the discretisation (g5.xlarge, $k_S a = 1$).}
\label{tab:precision}
\begin{tabular}{rr rr rrr}
\toprule
& & \multicolumn{2}{c}{c128 against Mie} & & \multicolumn{2}{c}{c64 against c128} \\
\cmidrule(lr){3-4}\cmidrule(lr){6-7}
$\nsub$ & cells & raw & volume-corrected & gmres\_tol & T-matrix & far field \\
\midrule
32 & 17\,256 & $4.97\times10^{-3}$ & $8.25\times10^{-4}$ & $10^{-5}$ & $4.8\times10^{-7}$ & $8.4\times10^{-7}$ \\
   &         &                     &                     & $10^{-6}$ & $6.4\times10^{-8}$ & $4.9\times10^{-9}$ \\
64 & 137\,376 & $7.14\times10^{-4}$ & $1.47\times10^{-4}$ & $10^{-5}$ & $5.1\times10^{-7}$ & $9.0\times10^{-7}$ \\
   &          &                     &                     & $10^{-6}$ & $8.1\times10^{-8}$ & $3.8\times10^{-9}$ \\
96 & 463\,400 & $2.69\times10^{-4}$ & $6.68\times10^{-5}$ & $10^{-5}$ & $5.1\times10^{-7}$ & $9.2\times10^{-7}$ \\
   &          &                     &                     & $10^{-6}$ & $8.6\times10^{-8}$ & $3.9\times10^{-9}$ \\
\bottomrule
\end{tabular}
\end{table}

The precision error is set by the GMRES tolerance, not by the word length. Tightening it from $10^{-5}$ to $10^{-6}$
lowers the far-field difference from about $9\times10^{-7}$ to $4\times10^{-9}$, below single precision's unit
roundoff, since the far field sums over all the cells; and it costs nothing measurable (31.3 against 30.8\,s at
$\nsub = 96$). At $10^{-6}$ the precision error is about 17\,000 times smaller than the discretisation error at
$\nsub = 96$, and the complex64 and complex128 far fields agree with Mie to the three figures shown at every size.

The discretisation error falls at about second order in $\nsub$ (the volume-corrected column: a factor 5.6 from 32 to
64 and 2.2 from 64 to 96, against 4 and 2.25 for second order), so a gmres\_tol of $10^{-5}$ would begin to matter
only at sizes far beyond these; $10^{-6}$ is the safe default. Against a target of $10^{-5}$ relative to the exact
solution, the limit at these sizes is the staircase, not the solver.

%=====================================================================
\section{Where the time goes}\label{sec:profile}

A profile of one solve at $\nsub = 96$ in complex128 on the g5.xlarge
(\code{scripts/data/profile_fft_sphere_gpu_g5.xlarge_complex128_n96.txt}) accounts for the stages the probe
did not separate (table~\ref{tab:profile}). GPU work is asynchronous, so the solves' time appears where the solver
waits for the device, inside \code{torch_gmres}.

\begin{table}[h]
\centering
\caption{One solve, $\nsub = 96$, complex128, g5.xlarge (52.3\,s in the solver), before the vectorisation below.}
\label{tab:profile}
\small
\begin{tabular}{lrl}
\toprule
Stage & seconds & where \\
\midrule
kernel build (\code{_point_canonical_kernel}) & 19.3 & CPU \\
\quad cell-averaged pair blocks (k-series) & 7.1 & \quad depends on the frequency \\
\quad symmetry and index tables & about 11 & \quad depends on $\nsub$ only \\
GMRES solves & 16.9 & GPU \\
incident patterns (\code{_build_incident_field_coupled}) & 11.7 & CPU, a loop over cells \\
copies from the device & 3.1 & \\
the solver's own lines, the composite T-matrix among them & 2.2 & CPU \\
symmetry applied to the solved columns & 1.2 & CPU \\
\bottomrule
\end{tabular}
\end{table}

The incident patterns were built cell by cell, an identity block and six small strain tensors for each of 463\,400
cells, though every entry is one product: 1 or $\tfrac12$ times a component of the cell's offset from the centre.
Built for all cells at once, they take 0.45\,s against 5.55\,s for the loop, on a two-core x86-64 machine (11.7\,s on
the g5.xlarge). The composite T-matrix, $\sum_n t_n \psi_n$ over the cells, is now formed in batches of 4096 cells
that stay in cache: 0.35 against 1.16\,s. Both are bit-identical to the loops, signed zeros included: NumPy sums
along the first axis in order, and the running sum enters each batch as its first row (the cheaper
$t\,\sum_n\psi_n$ differs at $10^{-15}$ and was not used). The four solvers that built the T-matrix by the loop all
use the new form, and their outputs, uniform and graded, in both precisions, are unchanged bit for bit.

With these, a complex64 solve at $\nsub = 96$ on the g5.xlarge should take roughly 26\,s instead of 37 (not yet
measured there), of which the kernel build is 18.5. That is the stage the design below moves off the GPU instance.

%=====================================================================
\section{A kernel stage and a solve stage}\label{sec:design}

The measurements separate the solver into two stages with different needs.

\begin{itemize}
\item The \emph{kernel stage} runs in complex128, on the CPU, and parallelises well (the eight Graviton cores ran
close to 800\% busy). More than half of it, at $\nsub = 96$, is symmetry and index tables that depend only on $\nsub$; the incident
patterns depend only on the cell centres. Only the pair blocks, 7.1\,s here, depend on the frequency and the medium.
\item The \emph{solve stage} runs in complex64, on the GPU, limited by memory bandwidth, and needs few CPU cores.
\end{itemize}

What passes between them is the canonical kernel, $K \times 9 \times 9$ complex numbers at the $K$ canonical
wavenumbers (table~\ref{tab:handover}), built in complex128 and rounded to complex64 for the solve.

\begin{table}[h]
\centering
\caption{The canonical kernel handed from the kernel stage to the solve stage.}
\label{tab:handover}
\begin{tabular}{rrrrr}
\toprule
$\nsub$ & $\nP$ & canonical wavenumbers $K$ & complex128 & complex64 \\
\midrule
64 & 135 & 54\,740 & 71\,MB & 35\,MB \\
96 & 225 & 246\,905 & 320\,MB & 160\,MB \\
128 & 315 & 669\,920 & 868\,MB & 434\,MB \\
\bottomrule
\end{tabular}
\end{table}

So the stages can run on different machines, with the kernel passed through object storage in seconds: the kernel
stage on many Graviton cores, as a batch array over frequencies on Spot instances, with the $\nsub$-only tables
computed once and kept; the solve stage on a small GPU instance, starting as each kernel is ready. In complex64 the
A10G's 24\,GB holds $\nsub = 128$ (19.6\,GB by extrapolation of the measured sizes); $\nsub = 160$ would need about
40\,GB, a 48\,GB GPU or fewer FFT temporaries (\S\ref{sec:memory}). Building this needs the solver split into a
function that builds and saves the canonical kernel and one that solves from it, with the present function their
composition; it has not been done.

%=====================================================================
\section{Limits of these measurements}\label{sec:limits}

\begin{itemize}
\item One frequency ($k_S a = 1$) and one moderate contrast, at which GMRES takes three or four steps. Stronger
contrast or higher frequency takes more, and the Krylov basis, now sized by the steps taken, grows with them.
\item The sphere is voxelised as a staircase; the discretisation error of table~\ref{tab:precision} is the
staircase's, and a body the lattice resolves smoothly would set a lower bar for the precision error to clear.
\item Each configuration was run once. Where two runs measured the same stage, they agreed to within about 6\% (the
product at $\nsub = 96$ on the CPU, 4.91 and 4.61\,s; the GPU solves, 16.97 and 16.85\,s), but no spread is
reported.
\item The memory attributed to FFT temporaries and to the k-series store is inferred from sizes, not profiled.
\item The timings of the vectorised stages (\S\ref{sec:profile}) were measured on a two-core x86-64 machine, not on
the instances.
\end{itemize}

%=====================================================================
\section{Reproduction}\label{sec:repro}

The data are in \code{scripts/data/}: \code{fft_sphere_scaling_r7g.2xlarge_cpu_complex128.jsonl},
\code{fft_sphere_scaling_g5.xlarge_cuda_complex128.jsonl}, \code{fft_sphere_scaling_g5.xlarge_cuda_complex64.jsonl},
\code{fft_sphere_precision_g5.xlarge_cuda.jsonl} and the profile above. The commits are 8818dab (the growing Krylov
basis), 4af0b78 (the data) and 9d72ca1 (the vectorised incident patterns and composite T-matrix). With the AWS
tooling of \code{aws/README.md}, from a checkout:

{\footnotesize
\begin{verbatim}
python3 aws/msc_aws.py run --on ec2 -- python scripts/probe_fft_sphere_scaling.py \
    --device cpu --n-sub 24,32,48,64,96 --warm-store
python3 aws/msc_aws.py run --on ec2 --gpu -- python scripts/probe_fft_sphere_scaling.py \
    --device cuda --dtype complex128 --n-sub 24,32,48,64,96 --warm-store
python3 aws/msc_aws.py run --on ec2 --gpu -- python scripts/probe_fft_sphere_scaling.py \
    --device cuda --dtype complex64 --tol 1e-6 --n-sub 24,32,48,64,96
python3 aws/msc_aws.py run --on ec2 --gpu -- python scripts/measure_fft_sphere_precision.py \
    --device cuda --n-sub 32,64,96
\end{verbatim}
}

\noindent
The same scripts run locally with \code{PYTHONPATH=.}; \code{--fit-only} reprints a table from its data file.
\code{--warm-store} grows the k-series store before the measured sizes, which otherwise include its one-time growth;
at $\nsub = 96$ that growth took roughly 45 minutes on the r7g.2xlarge.

\end{document}
